\section{A diagonal zero branch and infinitely many initial directions}
\label{diag:sec:main}

The fixed-index monotonicity theorem permits its derivative-order threshold
to depend on the Laurent index.  That dependence is essential.  We prove
that on the diagonal $k=n-1$, the smallest zero issuing from the elementary
endpoint moves initially to the right for infinitely many $n$ and initially
to the left for infinitely many $n$.  An exact inverse-function identity
organizes all these velocities.  The result concerns the central elementary
zero at $a=1$, which is separate from the branches near positive Bessel zeros
in the earlier adjacent-index expansion.

\subsection{The shifted lattice and its elementary endpoint}
\label{diag:subsec:lattice}

It is useful to allow an arbitrary lattice spacing $h>0$.  For integers
$n\geq0$, $k\geq1$, define
\begin{equation}\label{diag:eq:f}
 f_{n,k}(a)=(-1)^{n+k}\frac{d^k}{da^k}
                 \left(\frac{\log^n a}{a}\right),\qquad a>0,
\end{equation}
and
\begin{equation}\label{diag:eq:lattice}
 \mathscr F_{n,k}^{\rho,h}(a)
       =\sum_{m=0}^{\infty}\rho^m f_{n,k}(a+mh),
       \qquad 0\leq\rho\leq1.
\end{equation}
At $h=1$ this is the normalized Stieltjes--Lerch derivative family used in
the manuscript.  For $\rho<1$ it is $(-1)^{n+k}$ times the $k$th derivative
of the corresponding logarithmic Lerch series.  At $\rho=1$ we use
\eqref{diag:eq:lattice} itself; the undifferentiated logarithmic series need
not converge.  The derivative series converges absolutely because
\[
 f_{n,k}(x)=O_{n,k}\bigl(x^{-k-1}(1+\log^n x)\bigr)
 \quad(x\longrightarrow\infty).
\]
The convergence is locally uniform in $a>0$, together with any fixed
number of $a$-derivatives.  Near $\rho=0$ the series is jointly holomorphic
in $a$ and $\rho$, using the logarithm analytic near the positive axis.

The elementary polynomial is
\begin{equation}\label{diag:eq:P}
 P_{n,k}(x)=n![z^n]e^{xz}\prod_{j=1}^{k}\left(1+\frac zj\right).
\end{equation}
Differentiating $a^{-1-z}$ and comparing coefficients gives
\begin{equation}\label{diag:eq:elementary}
 f_{n,k}(a)=k!a^{-k-1}P_{n,k}(-\log a).
\end{equation}

\begin{lemma}[The central diagonal zero]\label{diag:lem:central}
For every $n\geq2$, $f_{n,n-1}$ has exactly $n$ simple positive zeros.
Its smallest zero is $1$, and
\begin{equation}\label{diag:eq:central-derivative}
 f_{n,n-1}'(1)=-n!.
\end{equation}
For every $h>0$ there is an $\epsilon_{n,h}>0$ such that
$\mathscr F_{n,n-1}^{\rho,h}$ has exactly $n$ simple positive zeros for
$0\leq\rho<\epsilon_{n,h}$.  Its smallest zero is a real analytic branch
$a_n^{(h)}(\rho)$ with $a_n^{(h)}(0)=1$.
\end{lemma}

\begin{proof}
The polynomial in \eqref{diag:eq:P} is obtained from $x^n$ by applying
$k$ operators $1+\partial_x/j$.  If a real-rooted polynomial $p$ has
distinct roots $r_i$ of multiplicities $m_i$, then away from these roots
the zeros of $p+cp'$, $c>0$, solve
\[
 \frac{p'(x)}{p(x)}=\sum_i\frac{m_i}{x-r_i}=-\frac1c.
\]
The left side is strictly decreasing on each complementary interval.
There is one new simple root between consecutive distinct roots and one
to the left of the smallest root.  An inherited root of multiplicity
$m_i\geq2$ has multiplicity $m_i-1$.  Starting from $x^n$, after $n-1$
operators all $n$ roots are simple; zero is the largest and the others
are negative.  Equation~\eqref{diag:eq:elementary} proves the first claim.

Since $n+(n-1)$ is odd,
\[
 f_{n,n-1}(a)=-\frac{d^{n-1}}{da^{n-1}}
                       \left(\frac{\log^n a}{a}\right).
\]
The expansion $\log^n a/a=(a-1)^n+O((a-1)^{n+1})$ proves
\eqref{diag:eq:central-derivative}.  The analytic implicit function theorem
therefore continues this zero, as well as each of the other simple zeros.

For completeness, no additional positive zeros enter from an endpoint.
For sufficiently large $a$, every term $f_{n,n-1}(a+mh)$ has the same
strict sign, because its argument exceeds the largest elementary zero.
As $a\downarrow0$, the $m=0$ term tends to $+\infty$, while the sum over
$m\geq1$ is bounded uniformly for $a$ in a small positive interval and
$0\leq\rho\leq1$.  Thus there are no zeros sufficiently close to zero.
On the remaining compact set outside small neighborhoods of the elementary
zeros, uniform convergence to $f_{n,n-1}$ as $\rho\downarrow0$ excludes
zeros.  Inside each neighborhood, uniform convergence of the first
$a$-derivative preserves its strict sign and hence gives exactly one
simple zero.  The branch issuing from $1$ is consequently the smallest.
\end{proof}

\subsection{An inverse function generates every diagonal velocity}
\label{diag:subsec:inverse}

Let $\stirling{n}{j}$ denote the unsigned Stirling numbers of the
first kind, fixed by
\[
 \prod_{r=0}^{n-1}(z+r)=\sum_{j=0}^n\stirling{n}{j}z^j.
\]

\begin{theorem}[Exact velocity identity]\label{diag:thm:generator}
For $A>1$ and $n\geq1$, put
\begin{align}
 v_n(A)
  &=-\frac1{n!}\left.\frac{d^{n-1}}{dx^{n-1}}
                  \left(\frac{\log^n x}{x}\right)\right|_{x=A}
                  \label{diag:eq:velocity-derivative}\\
  &=A^{-n}\sum_{j=1}^n
       \stirling{n}{n-j+1}\frac{(-\log A)^j}{j!}.
                  \label{diag:eq:velocity-stirling}
\end{align}
For the lattice spacing $h=A-1$, the smallest diagonal branch satisfies
\begin{equation}\label{diag:eq:root-velocity}
 \bigl(a_n^{(h)}\bigr)'(0)=v_n(A),\qquad n\geq2.
\end{equation}
There is a unique analytic germ $u_A(t)$ at zero satisfying
\begin{equation}\label{diag:eq:inverse-equation}
 e^{u_A(t)}=A+t u_A(t),\qquad u_A(0)=\log A.
\end{equation}
Its coefficients give the convergent ordinary generating identity
\begin{equation}\label{diag:eq:velocity-generator}
 V_A(t):=\sum_{n=1}^{\infty}v_n(A)t^n=\log A-u_A(t)
\end{equation}
for all sufficiently small complex $t$.
\end{theorem}

\begin{proof}
Differentiate the zero equation at $\rho=0$.  The shifted series and
\eqref{diag:eq:central-derivative} give
\[
 \bigl(a_n^{(h)}\bigr)'(0)
   =-\frac{f_{n,n-1}(1+h)}{f_{n,n-1}'(1)}
   =\frac{f_{n,n-1}(A)}{n!},
\]
which is \eqref{diag:eq:velocity-derivative}.  The Stirling identity implies
\[
 P_{n,n-1}(x)
       =n\sum_{j=1}^n\stirling{n}{n-j+1}\frac{x^j}{j!}.
\]
Together with \eqref{diag:eq:elementary}, this proves
\eqref{diag:eq:velocity-stirling}.  The same finite identity holds at $n=1$
by direct calculation; that term is included for the generating function.

Let $y(t)$ solve $y=A+t\log y$, with $y(0)=A$.  The analytic implicit
function theorem supplies a unique germ.  The Lagrange--B\"urmann formula
with outer function $\log y$ gives
\[
 [t^n]\log y(t)
   =\frac1{n!}\left.\frac{d^{n-1}}{dx^{n-1}}
                        \left(\frac{\log^n x}{x}\right)\right|_{x=A}.
\]
This is the standard analytic form of Lagrange inversion
\cite{Gessel2016}.  Set $u_A=\log y$.  Its equation is
\eqref{diag:eq:inverse-equation}, and coefficient comparison proves
\eqref{diag:eq:velocity-generator}.
\end{proof}

\begin{corollary}[An exact quadratic recurrence]\label{diag:cor:recurrence}
Write $L=\log A$ and $p_n(L)=A^n v_n(A)$.  Then $p_1(L)=-L$, and for
$n\geq1$,
\begin{equation}\label{diag:eq:recurrence}
 (n+1)p_{n+1}(L)
   =(n+1-nL)p_n(L)
       +\frac n2\sum_{j=1}^{n-1}p_j(L)p_{n-j}(L).
\end{equation}
Each $p_n$ is a rational polynomial of degree $n$, with constant term zero,
linear coefficient $-1$, and leading coefficient $(-1)^n/n$.
\end{corollary}

\begin{proof}
The generating identity is equivalent to
\[
 A e^{-V_A(t)}=A+t(L-V_A(t)).
\]
Differentiation and elimination of the exponential yield
\begin{equation}\label{diag:eq:ode}
 \bigl(A+t(L-1-V_A(t))\bigr)V_A'(t)=V_A(t)-L.
\end{equation}
The constant coefficient gives $v_1=-L/A$.  At degree $n\geq1$, the
convolution is
\[
 \sum_{j=1}^{n-1}(n-j)v_jv_{n-j}
       =\frac n2\sum_{j=1}^{n-1}v_jv_{n-j}.
\]
Multiplication by $A^n$ gives \eqref{diag:eq:recurrence}.  The coefficient
statements follow directly from \eqref{diag:eq:velocity-stirling}, using
$\stirling{n}{n}=1$ and
$\stirling{n}{1}=(n-1)!$.
\end{proof}

\begin{corollary}[Nonvanishing at algebraic lattice endpoints]
\label{diag:cor:algebraic-nonvanishing}
If $A>1$ is algebraic, then $v_n(A)\ne0$ for every $n\ge1$.
\end{corollary}
\begin{proof}
The Hermite--Lindemann theorem implies that $\log A$ is transcendental:
it is nonzero, and its exponential is algebraic. Each $p_n$ in
Corollary~\ref{diag:cor:recurrence} is a nonzero rational polynomial.
Thus $p_n(\log A)\ne0$, and $v_n(A)=A^{-n}p_n(\log A)\ne0$.
\end{proof}

\subsection{A sign theorem without an unproved singularity analysis}
\label{diag:subsec:signs}

\begin{lemma}[Positive-axis continuation and coefficient signs]
\label{diag:lem:pringsheim}
Let $g$ be a nonpolynomial analytic germ at zero with real Taylor
coefficients.  Suppose it admits coherent holomorphic continuation to a
neighborhood of every point of $[0,\infty)$ and has at most polynomial
growth on the positive real axis.  Then its Taylor coefficients include
infinitely many positive and infinitely many negative terms.
\end{lemma}

\begin{proof}
Suppose the coefficients have one sign from some index $N$ onward.
Subtract the Taylor polynomial of degree $N-1$ and, if needed, multiply
by $-1$.  We obtain a nonpolynomial germ
\[
 H(t)=\sum_{n\geq N}b_n t^n,\qquad b_n\geq0,
\]
with the same continuation and growth properties.

Here is the needed positive-coefficient argument in full.  If the radius
of convergence were a finite $R>0$, analyticity near the positive point
$R$ would imply, for every $j\geq0$,
\[
 H^{(j)}(R)=\sum_{n\geq j}(n)_j b_n R^{n-j}\geq0.
\]
Indeed the equality holds first at $0<t<R$; monotone convergence of its
nonnegative terms and continuity of $H^{(j)}$ at $R$ give the displayed
formula.  Choose $\delta>0$ within the Taylor radius at $R$.  That Taylor
series converges absolutely, so Tonelli's theorem and the binomial theorem
give
\[
 \sum_{j\geq0}\frac{H^{(j)}(R)}{j!}\delta^j
     =\sum_{n\geq0}b_n(R+\delta)^n<\infty.
\]
The original power series would converge beyond $R$, a contradiction.
This is the usual Pringsheim argument, and it proves that $H$ is entire.

Now $0\leq b_n t^n\leq H(t)$ for every positive real $t$.  A polynomial
upper bound $H(t)=O(t^M)$ forces $b_n=0$ for every $n>M$.  Thus $H$, and
hence $g$, is a polynomial, the final contradiction.  Eventual
nonnegativity and eventual nonpositivity have both been excluded.
\end{proof}

\begin{theorem}[Infinite diagonal oscillation]\label{diag:thm:oscillation}
For every $A>1$, each of the sets
\[
 \{n\geq1:v_n(A)>0\},\qquad \{n\geq1:v_n(A)<0\}
\]
is infinite.  Consequently, for every lattice spacing $h>0$, the smallest
zero branch of $\mathscr F_{n,n-1}^{\rho,h}$ initially moves in each
direction for infinitely many Laurent indices $n\geq2$.
\end{theorem}

\begin{proof}
Set $L=\log A>0$ and consider
\[
 h_A(u)=\frac{e^u-A}{u},\qquad u\geq L.
\]
It has $h_A(L)=0$ and tends to infinity.  Its derivative is
\[
 h_A'(u)=\frac{e^u(u-1)+A}{u^2}>0.
\]
For the strict inequality, the numerator equals $AL>0$ at $u=L$, and its
derivative is $ue^u>0$ for $u>0$.  Thus $h_A$ has a unique increasing real
inverse on $[0,\infty)$.  The derivative never vanishes, so the complex
implicit function theorem continues this inverse holomorphically near
every positive real point.  It is the continuation of $u_A$ from
\eqref{diag:eq:inverse-equation}.

The inverse is unbounded but has logarithmic growth.  To see the latter
directly, take $U=2\log(t+A+2)$.  For all sufficiently large $t$,
\[
 e^U=(t+A+2)^2>tU+A,
\]
so $h_A(U)>t$ and monotonicity implies
$u_A(t)<2\log(t+A+2)$.  An unbounded function of this growth cannot be a
polynomial.  Lemma~\ref{diag:lem:pringsheim} applies to $u_A$ and proves
infinitely many Taylor coefficients of each sign.  Equation
\eqref{diag:eq:velocity-generator} negates all its nonconstant
coefficients.  This proves the result for $v_n(A)$; omission of the single
index $n=1$ does not change infinitude.
\end{proof}

\begin{corollary}[A restriction on uniform monotonicity thresholds]
\label{diag:cor:threshold}
For the unit lattice, no integer $N$ has the property that every smallest
diagonal branch $a_n^{(1)}(\rho)$, $n\geq N$, is nonincreasing near
$\rho=0$.  If $K_n$ is any threshold guaranteeing decreasing motion of
every zero branch throughout the Lerch interval for every $k\geq K_n$,
then $K_n>n-1$ for infinitely many $n$.
\end{corollary}

\begin{proof}
For infinitely many $n$, Theorem~\ref{diag:thm:oscillation} gives
$\bigl(a_n^{(1)}\bigr)'(0)>0$.  Analyticity makes the branch strictly
increasing on a sufficiently small positive interval, contradicting each
of the proposed conclusions when $k=n-1$ is included.
\end{proof}

This conclusion respects the quantifiers of the earlier fixed-index
theorem.  It does not assert that its threshold fails to exist, or that
the diagonal branch must continue simply throughout $0\leq\rho\leq1$.
The local interval in Lemma~\ref{diag:lem:central} can depend on $n$.

\subsection{A short exact example and the verification artifacts}
\label{diag:subsec:verification}

For the unit lattice put $L=\log2$.  The first relevant velocities are
\begin{align*}
 v_2(2)&=\frac{L^2-2L}{8},\\
 v_3(2)&=\frac{-2L^3+9L^2-6L}{48},\\
 v_4(2)&=\frac{3L^4-22L^3+36L^2-12L}{192}.
\end{align*}
The fourth-index example is positive:
\begin{equation}\label{diag:eq:positive-certificate}
 0.01221096503<v_4(2)<0.01221096504.
\end{equation}
For a reproducible exact certificate, use
\[
 L_M=2\sum_{j=0}^{M-1}\frac1{(2j+1)3^{2j+1}},\qquad
 0<\log2-L_M<
     \frac{2}{(2M+1)3^{2M+1}(1-1/9)}.
\]
This is the positive series for $2\operatorname{arctanh}(1/3)$ with a
geometric tail majorant.  At $M=20$, rational interval Horner evaluation
of the displayed polynomial gives \eqref{diag:eq:positive-certificate}.
Equivalently,
\[
 P_{4,3}(x)=x^4+\frac{22}{3}x^3+12x^2+4x,
 \qquad v_4(2)=\frac{P_{4,3}(-\log2)}{64}.
\]

The exact verifier compares the Stirling expression with the quadratic
recurrence over $\mathbb Q[L]$ through $n=40$.  It also checks the inverse
generating equation through degree forty at a formal rational value of
$L$, and encloses the actual velocities at $L=\log2$ by rational
arithmetic.  The signs for $2\leq n\leq12$ are
\[
\begin{array}{c|rrrrrrrrrrr}
 n&2&3&4&5&6&7&8&9&10&11&12\\ \hline
 \operatorname{sgn}v_n(2)&-&-&+&+&+&-&-&+&+&+&-
\end{array}
\]
These finite certificates check the implementation and the short example.
The infinite sign conclusion rests on the analytic proof above.

\subsection{Questions left by the inverse representation}
\label{diag:subsec:questions}

The exact generating function permits a quantitative refinement, but its
complex branches must be handled carefully.  A finite critical point of
the inverse equation satisfies
\begin{equation}\label{diag:eq:critical-values}
 (1-u)e^u=A,\qquad t=e^u.
\end{equation}
Different Lambert-$W$ sheets give different critical values, some tending
to zero on other sheets.  Therefore one cannot read off the radius of
the germ $u_A(t)$ by minimizing over all formal critical values.  The
oscillation theorem does not require such a dominance assertion.

Concrete further questions are to identify the accessible dominant
singularities and obtain asymptotic formulas for $v_n(A)$; to bound gaps
between indices with opposite velocity signs; to determine a uniform
neighborhood in which the first velocity controls the true branch; and
to extend the inverse-series method to $n=k+r$, $r>1$, where the central
elementary root has multiplicity $r$ and its unfolding is generally
described by fractional powers.  These questions connect the central
branch treated here with the existing positive-Bessel-zero regime.
